\documentclass[UTF8,fontset=none,11pt,a4paper]{ctexart}
\usepackage{ipara-handbook}
\usepackage{amsmath,booktabs,tabularx,hyperref}
\handbooksetup{OS-B02 Process × Virtual Memory}{408 · OS Internal Bridge}{Identity · Address Space · COW · Fault}
\begin{document}
\handbooktitle{Process × Virtual Memory}{进程身份怎样绑定地址空间，并在 fork 后延迟分化}{408 · OS-B02}
\begin{corebox}{Mother Interface}
\[
\boxed{\text{Process Identity}\to\text{Address-Space Association}\to\text{Shared Mapping}\to\text{Write Fault}\to\text{COW Divergence}}
\]
本 Bridge 连接“谁在执行”和“它看到哪套地址空间”，不重新定义进程生命周期或页表硬件。
\end{corebox}
\begin{methodbox}{Owners 与边界}
OS01 Owns process identity、fork/exec 与调度；OS04 Owns address space、PTE、page fault、frame 与 COW 策略。本册 Owns 二者的绑定、共享到分化的交接；硬件 VA\(\to\)PA 进入 X-B02。
\end{methodbox}
\section{共享不是复制}
fork 后父子进程通常共享物理页，但各自拥有可区分的地址空间元数据；可写页被暂时标为只读并增加引用关系。读访问仍共享，任一方写入时才触发 fault，OS 分配新页、复制内容、更新写入方映射，随后重试写指令。
\section{最小状态轨迹}
\[
\text{parent}\xrightarrow{fork}\text{parent+child, shared page}\xrightarrow{write}\text{fault}\xrightarrow{copy}\text{private writable page}.
\]
进程切换时，切换的是“当前执行实体 + 地址空间关联”；不能把 context switch 自动等同于复制全部页。
\section{最小例子与验证}
父进程写入变量前，子进程读取同一变量。若先读后写，读到相同内容；写入方触发 COW 后，父子后续读值可以不同。题目问“是否共享”时必须区分共享虚拟地址、共享物理页、共享文件对象和共享后续写入结果。
\begin{warnbox}{Anti-Bridge}
fork 的共享映射 $\neq$ 永久共享可写物理页；COW fault $\neq$ 普通 TLB miss；进程身份 $\neq$ 一组固定物理地址。
\end{warnbox}
\begin{corebox}{Compression}
当前进程是谁？它绑定哪套 address space？页目前共享什么？哪次写入触发分化？修复后谁的映射被改？
\end{corebox}
\end{document}
